\documentclass{article}
\usepackage{amsmath,amssymb,amsthm}
\begin{document}

The space \(H^1_0(\Omega)\) is a Hilbert space with the \(H^1\) norm \(\|\cdot\|_{1,\Omega}\). Since \(F\in H^{-1}(\Omega)\), it is a continuous linear functional on this space. The assumed boundedness and coercivity of \(a\) therefore allow the Lax--Milgram lemma to be applied, giving a unique \(y\in H^1_0(\Omega)\) such that \(a(y,v)=F(v)\) for every \(v\in H^1_0(\Omega)\).

The finite-dimensional subspace \(V_h^0\), equipped with the inherited norm, is also a Hilbert space. The restrictions of \(a\) and \(F\) to \(V_h^0\) remain bounded, and the restriction of \(a\) remains coercive with the same constant \(c\). A second application of the Lax--Milgram lemma thus gives a unique \(y_h\in V_h^0\) satisfying \(a(y_h,v_h)=F(v_h)\) for every \(v_h\in V_h^0\).

Set \(e=y-y_h\). Subtracting the two variational equations for any \(v_h\in V_h^0\) yields Galerkin orthogonality: \(a(e,v_h)=0\). Now fix an arbitrary \(v_h\in V_h^0\). Because \(y_h-v_h\in V_h^0\), orthogonality, coercivity, and boundedness give
\[
c\|e\|_{1,\Omega}^2\leq a(e,e)=a(e,y-v_h)\leq C\|e\|_{1,\Omega}\|y-v_h\|_{1,\Omega}.
\]
If \(e\neq0\), division by \(\|e\|_{1,\Omega}\) gives \(\|e\|_{1,\Omega}\leq(C/c)\|y-v_h\|_{1,\Omega}\); the same inequality is immediate if \(e=0\). Since \(v_h\) was arbitrary, taking the infimum over \(V_h^0\) proves
\[
\|y-y_h\|_{1,\Omega}\leq\frac{C}{c}\inf_{v_h\in V_h^0}\|y-v_h\|_{1,\Omega}.
\]

\end{document}